\section{上下文管理：给 Agent 一张地图}

上下文管理是让 Agent 高效处理大型复杂任务的最大挑战之一。OpenAI 团队在这方面积累了深刻的教训。

\subsection{``一个大文件''方法的失败}

团队最初尝试将所有指导信息放入一个大型 \texttt{AGENTS.md} 文件中。这种方法以可预见的方式失败了：

\begin{warningbox}{巨型 AGENTS.md 的四大问题}
\begin{enumerate}[nosep]
  \item \textbf{上下文是稀缺资源}：一个巨型指令文件会挤占任务本身、代码和相关文档的空间，导致 Agent 遗漏关键约束或优化错误的目标。
  \item \textbf{过多指导等于没有指导}：当所有东西都是``重要''的，就没有东西是重要的。Agent 最终只会进行局部模式匹配，而非有意识地导航。
  \item \textbf{瞬间腐烂}：一个单体手册很快变成过时规则的坟场。Agent 无法判断哪些仍然有效，人类也懒得维护，文件悄然变成一个``有吸引力的危险物''（attractive nuisance）。
  \item \textbf{难以验证}：单一大文件不利于机械化检查（覆盖率、新鲜度、归属、交叉链接），因此漂移不可避免。
\end{enumerate}
\end{warningbox}

\subsection{地图而非百科全书}

团队转而采用\textbf{``目录 + 深层知识库''}的架构：

\begin{importantbox}{AGENTS.md 作为目录（Table of Contents）}
将 \texttt{AGENTS.md} 视为\textbf{目录}而非百科全书。它应当是简短的（约 100 行），仅包含指向更深层信息源的指针。仓库的知识库存放在结构化的 \texttt{docs/} 目录中，作为``唯一真实来源''（system of record）。
\end{importantbox}

文章给出的仓库知识架构如下：

\begin{table}[H]
\centering
\small
\begin{tabular}{p{5cm}p{8cm}}
\toprule
\textbf{文件/目录} & \textbf{用途} \\
\midrule
\texttt{AGENTS.md} & 约 100 行的入口地图，指向深层文档 \\
\texttt{ARCHITECTURE.md} & 顶层架构图，描述域与包的层次关系 \\
\texttt{docs/design-docs/} & 设计文档索引，含验证状态和核心信念 \\
\texttt{docs/exec-plans/} & 执行计划（活跃/已完成），含决策日志 \\
\texttt{docs/generated/} & 自动生成的文档（如数据库 schema） \\
\texttt{docs/product-specs/} & 产品规格说明书 \\
\texttt{docs/references/} & 外部参考资料的 LLM 友好版本 \\
\texttt{QUALITY\_SCORE.md} & 各产品域和架构层的质量评分 \\
\texttt{RELIABILITY.md} & 可靠性要求 \\
\texttt{SECURITY.md} & 安全规范 \\
\bottomrule
\end{tabular}
\caption{知识库目录结构及用途}
\end{table}

\subsection{渐进式披露（Progressive Disclosure）}

这种架构实现了\textbf{渐进式披露}：Agent 从一个小而稳定的入口点开始，被教导在需要时去哪里查找更多信息，而不是一开始就被信息淹没。

\begin{knowledgebox}{渐进式披露的设计理念}
渐进式披露（Progressive Disclosure）是交互设计中的经典原则，指的是只在用户需要时才展示详细信息。在 Agent 上下文管理中，这意味着：Agent 的初始上下文只包含``地图''级别的信息，具体细节由 Agent 在任务执行过程中按需获取。这与人类新员工入职的逻辑相同——先给组织架构图，再按需深入各个模块。
\end{knowledgebox}

\subsection{机械化的知识库维护}

团队不仅依赖人工维护知识库，还通过\textbf{机械化手段}确保其质量：

\begin{itemize}[nosep]
  \item 专用的 Linter 和 CI Job 验证知识库是否最新、交叉链接正确、结构合规
  \item 定期运行的``文档园丁''（doc-gardening）Agent 扫描过时或不反映真实代码行为的文档
  \item 发现问题后自动开启修复性 Pull Request
\end{itemize}

\subsection{本章小结}

上下文管理的核心原则是``给 Agent 一张地图，而非一本百科全书''。具体实践包括：将 AGENTS.md 保持简短作为入口指针；将详细知识存入结构化的 docs/ 目录；采用渐进式披露策略；并通过 Linter、CI 和自动化 Agent 机械化地维护知识库的新鲜度和一致性。

\newpage

